% !TEX TS-program = lualatex

%%%
% src::
%   url = https://github.com/pfradin/luadraw/discussions/96#discussioncomment-14477291
%%%

\documentclass[border = 3pt]{standalone}

%%%
% Des couleurs faciles d'emploi via le package \xcolor!
%%%
\usepackage[svgnames]{xcolor}

%%%
% La bibliothèque ''luadraw'' allie une facilité d’utilisation
% à un rendu particulièrement soigné.
%%%
\usepackage[3d]{luadraw}

\begin{document}

\begin{luadraw*}{name = box-falling-walls}
local ld = luadraw

local Origin = ld.pt3d.Origin
local M      = ld.pt3d.M
local vecI   = ld.pt3d.vecI
local vecJ   = ld.pt3d.vecJ
local vecK   = ld.pt3d.vecK

------
-- Préréglages.
------
local a, b, h = 4, 3, 5

local shiftbox = 3*b

------
-- ¨Def de la zone graphique.
------
local xmin, xmax = -5*a, a
local ymin, ymax = -b - shiftbox, b
local zmin, zmax = -h, h

local graphview = ld.graph3d:new{
  window3d = {xmin, xmax, ymin, ymax, -h, h},
  viewdir  = {195, 70},
  adjust2d = true,
  bbox     = false,
}

------
-- Réglages liés aux lignes.
------
ld.Hiddenlines     = true
ld.Hiddenlinestyle = "dashed"

------
-- Les ¨elts géométriques de la boîte.
------
local box = ld.parallelep(
  Origin,
  a*vecI, b*vecJ, h*vecK
)

local A, B, C, D, E, F, G, H = table.unpack(box.vertices)

local box_facets = ld.poly2facet(box)

------
-- Déplacements des faces en retirant celle non voulue.
------
box_facets[4] = ld.rotate3d(box_facets[4], -30, {C, C-D})
box_facets[3] = ld.rotate3d(box_facets[3], 30 , {A, B-A})
box_facets[5] = ld.rotate3d(box_facets[5], -30, {A, D-A})
box_facets[6] = ld.rotate3d(box_facets[6], 30 , {B, C-B})

table.remove(box_facets, 2)

------
-- La boîte ouverte.
------
local box = ld.shift3d(box, -shiftbox*vecJ)

table.remove(box.facets, 2)

------
-- Tracé des ¨elts graphiques.
------
graphview:Dscene3d(
  graphview:addFacet(
    box_facets,
    {
      color     = "LightGray",
      contrast  = 0.5,
      edge      = true,
      edgecolor = "gray"
    }
  )
)

graphview:Dpoly(
  box,
  {mode = ld.mShadedHidden}
)

------
-- Montrons le résultat de notre oeuvre.
------
graphview:Show()
\end{luadraw*}

\end{document}
